% Single-file draft paper for supervisor review — not the thesis template.
% Prognosis/survival (PROGNOSER) work is intentionally left out per request.
\documentclass[11pt,a4paper]{article}

\usepackage[utf8]{inputenc}
\usepackage[T1]{fontenc}
\usepackage[margin=2.5cm]{geometry}
\usepackage{graphicx}
\usepackage{amsmath,amssymb}
\usepackage{booktabs}
\usepackage{microtype}
\usepackage{hyperref}
\hypersetup{colorlinks=true, linkcolor=blue, citecolor=blue, urlcolor=blue}

\sloppy
\emergencystretch=3em

% Inline draft/caveat callout — flags content that is not yet settled.
\newcommand{\draftnote}[1]{%
  \par\noindent\textbf{Draft note:} \textit{#1}\par\vspace{0.5\baselineskip}%
}

\title{Longitudinal Graph Representation Learning\\
  for Early Detection of Alzheimer's Disease\\[0.3em]
  \large Draft --- Progress Summary for Supervisor Review}
\author{Fraj Yassine Lakhal}
\date{\today}

\begin{document}
\maketitle

\begin{abstract}
We study early detection of Alzheimer's disease progression from longitudinal
resting-state fMRI in the DELCODE cohort. A graph attention/variational
autoencoder (GAAE/VGAE) is pretrained on static functional-connectivity graphs
to produce frozen per-visit embeddings; a recurrent trajectory head
(GE-LSTM/GE-GRU), conditioned on inter-scan time elapsed, consumes the
resulting embedding sequence to predict MCI-to-AD conversion. Three findings
have emerged that reshape the argument. First, a published
spatiotemporal-transformer baseline (Brain-TokenGT) ships with its recurrent
core untrained, diverges when it is enabled, and after stabilisation is
\emph{not reproducible at a fixed seed}: across 19 runs of one configuration
within-seed SD ($0.1011$) exceeds between-seed SD ($0.0706$), while our own
pipeline re-runs byte-identically. Second, a pre-registered four-arm ablation
finds that removing the graph encoder entirely does not degrade performance,
extending to the longitudinal small-$N$ regime a result previously shown only
cross-sectionally at large $N$. Third, $\Delta t$ conditioning cannot be
evidenced on this cohort at all, because 90\% of DELCODE's inter-visit
intervals are identical by protocol. This draft omits the
time-to-event/survival extension (PROGNOSER), which is scoped rather than
delivered.
\end{abstract}

\section{Introduction \& Motivation}
\begin{itemize}
  \item The AD continuum: Subjective Cognitive Decline (SCD) $\to$ Mild
    Cognitive Impairment (MCI) $\to$ AD dementia, and the push toward early
    biomarker confirmation now that disease-modifying therapies (Lecanemab /
    Donanemab) require it.
  \item Resting-state fMRI as a non-invasive window onto functional
    dysconnectivity that can precede macrostructural atrophy, versus the
    cost/invasiveness of PET and CSF gold standards.
  \item Core modeling challenge: high-dimensional graph data under extreme
    sample constraints ($N < 150$), with irregular longitudinal follow-up
    ($\Delta t$) and patient attrition.
  \item Contributions: (1) self-supervised brain-graph pretraining on static
    connectomes, (2) time-aware recurrent modeling of longitudinal functional
    trajectories, (3) an empirical audit and repair of a published SOTA
    baseline.
\end{itemize}

\section{Related Work}
\begin{itemize}
  \item Brain connectivity \& parcellation: BOLD-derived FC matrices;
    Schaefer 100/200, Harvard-Oxford, AAL parcellations; Yeo functional
    networks (DMN, Limbic, Dorsal Attention).
  \item Graph representation learning: GCN / GAT / GATv2 message passing;
    self-supervised graph autoencoders (GAE/VGAE) as pretext tasks.
  \item Longitudinal \& spatiotemporal modeling: LSTM/GRU with continuous
    $\Delta t$ integration; spatiotemporal graph transformers and dynamic
    graph updates (Brain-TokenGT, EvolveGCN).
  \item Gaps this work targets: cross-sectional evaluation bias in prior
    work, under-audited pretext-to-downstream leakage, and unexamined
    failure modes of high-capacity transformers on small clinical cohorts.
\end{itemize}

\section{Data \& Preprocessing}
\begin{itemize}
  \item DELCODE cohort: diagnostic groups (\texttt{prmdiag}), longitudinal
    follow-up structure, Converter-vs-Stable-MCI definition, scanning
    interval statistics.
  \item External validation on ADNI and OASIS-3 (harmonization, visit
    alignment, cohort matching).
  \item fMRI preprocessing: motion correction, nuisance regression, bandpass
    filtering, linear-model age-effect regression.
  \item Leakage prevention: nested patient-level stratified cross-validation
    with strictly isolated held-out test manifests; historical
    pretext-to-downstream leakage identified and closed.
\end{itemize}

\section{Methodology}

\subsection{Architecture Overview}
Raw BOLD $\to$ FC graphs $\to$ frozen latent embedding (Stage 1) $\to$
recurrent trajectory classifier (Stage 2) $\to$ conversion risk output.

\subsection{Stage 1 --- Static Graph Pretraining}
\begin{itemize}
  \item Graph Attention Autoencoder (GAAE) and Variational Graph Autoencoder
    (VGAE); GATv2 message passing, cosine-similarity reconstruction, latent
    mean pooling.
  \item Latent dimensionality selection via latent-space analysis and FDR
    feature filtering.
\end{itemize}

\subsection{Stage 2 --- Longitudinal Trajectory Modeling (GE-LSTM / GE-GRU)}
\begin{itemize}
  \item Recurrent gate updates conditioned on inter-scan time elapsed
    ($\Delta t$).
  \item LayerNorm, dropout, and frozen-encoder BatchNorm kept in eval mode.
\end{itemize}

\subsection{Baselines \& Optimization}
\begin{itemize}
  \item Static baselines: Graph Encoder Classifier (GEC), Graph Embedding
    Perceptron (GEP), direct raw-FC classifier.
  \item Class imbalance handled via weighted cost averaging; decision
    threshold selected on validation via $F_1$-optimal or Youden's $J$, never
    from test metrics.
\end{itemize}

\section{Baseline Audit: Brain-TokenGT}
\begin{itemize}
  \item Brain-TokenGT tokenizes ROI cliques and applies cross-visit
    self-attention; evaluated here as a published spatiotemporal-transformer
    benchmark.
  \item In the as-released implementation, the recurrent core is disabled
    (\texttt{train\_give=False}), and end-to-end training of the
    EvolveGCN-H/GRCU component is unstable (gradient explosions, NaNs)
    without further intervention.
  \item Partial repair applied: parameter-group decoupling, learning-rate
    damping ($0.1\times$) and weight decay ($10^{-3}$) on the recurrent
    parameters, which reduces outright training collapse.
\end{itemize}

\draftnote{\textbf{Resolved, 2026-08-22 --- the audit is complete and the
  statistics previously reported here are withdrawn.} Nineteen completed runs of the
  identical stabilised configuration (7 at seed 42, 4 each at seeds 43/44/45) give a grand
  mean test AUC of $0.6025$ (pooled SD $0.1123$), against the $0.6672 \pm 0.1008$
  previously quoted from one hand-picked run per seed. The variance decomposition is the
  finding: mean within-seed SD $= 0.1011$ exceeds between-seed-mean SD $= 0.0706$, at every
  seed and not only at 42, where the range spans $0.3571$--$0.7078$. ``Stabilised''
  therefore means \emph{crashes less often}, not \emph{reproducible}. Both paired
  significance tests are consequently invalid --- the $n{=}4$ test-AUC $t$-test
  ($p = 0.0299$) treats seed as the only variance source when within-seed variance is
  larger, and the $n{=}20$ CV-AUC $t$-test ($p = 1.59\times10^{-6}$) treats correlated
  folds as independent --- and are dropped rather than corrected. The \emph{direction}
  survives and should be reported as two distributions plus an effect size: GE-LSTM
  $0.8782 \pm 0.0256$ over 4 seeds, against a Brain-TokenGT distribution whose best run
  ever observed ($0.8182$) sits below the GE-LSTM mean. Mechanism: the repo seeds
  \texttt{cudnn.deterministic} but never calls
  \texttt{torch.use\_deterministic\_algorithms(True)} nor sets
  \texttt{CUBLAS\_WORKSPACE\_CONFIG}, so \texttt{TopKPooling} and an index-gather in the
  GRCU backward hit nondeterministic CUDA kernels. Our \texttt{none}-arm GE-LSTM, which
  executes no scatter backward, re-runs byte-identically --- making
  \emph{``our pipeline is reproducible and theirs is not''} the stronger claim.}

\begin{itemize}
  \item Fairness caveats: cohort-windowing disparity between the baseline's
    published setup (1--6 visits) and our matched-cohort setup (2--3 visits);
    parameter-count disparity ($\sim$1{,}380 tokens / 604k params for
    Brain-TokenGT vs.\ 31k--520k params for the recurrent trajectory models).
\end{itemize}

\section{Results \& Ablations}

\subsection{Longitudinal Conversion Prediction}
5-fold cross-validation and held-out test evaluation on DELCODE, reported as
AUC-ROC, sensitivity, specificity, and balanced accuracy.
\draftnote{Insert current best-model table here once the reconstruction-value
  ablation branch (\texttt{ablation/reconstruction-value}) settles.}

\subsection{Architectural Ablation Suite}
\begin{itemize}
  \item Encoder (4 pre-registered arms, 4 seeds): \texttt{none} (no encoder,
    31k params) $0.8313 \pm 0.0529$; \texttt{pretrained\_frozen} (520k)
    $0.7812 \pm 0.0122$; \texttt{random} $0.7312 \pm 0.1113$;
    \texttt{pretrained\_finetuned} $0.7107 \pm 0.1946$. Metadata-only floor:
    CV AUC $0.6157 \pm 0.0653$, test AUC $0.4929$.
  \item Temporal head: LSTM vs.\ GRU vs.\ static pooling.
  \item Trajectory ablations (eval-time, fixed checkpoint): shuffling visit order costs
    $0.9071 \rightarrow 0.8486 \pm 0.0258$; zeroing $\Delta t$ costs \emph{nothing}
    ($0.9071 \rightarrow 0.9071$, identical sens/spec).
  \item \textbf{Parcellation (whole brain vs.\ DMN vs.\ DMN + limbic): cut.} No such
    experiment exists in the registry; it moves to future work rather than being implied.
\end{itemize}
\draftnote{\textbf{The fine-tuning conclusion is withdrawn, not caveated.} The
  \texttt{pretrained\_finetuned} arm optimises the unfrozen pretrained GATv2 encoder and
  the freshly initialised head at a \emph{single shared} learning rate, so it measures
  naive fine-tuning rather than fine-tuning. It is additionally confirmed
  non-deterministic (a re-run at seed 43 diverged on folds 1 and 4), so its $\pm 0.1946$
  conflates seed and run noise exactly as Brain-TokenGT's does. A repeat sweep
  ($R \ge 3 \times 4$ seeds) is running; an \texttt{encoder\_lr\_scale} arm is needed
  before any claim about fine-tuning can be made.}

\subsection{Early Detection Horizon}
Discriminative capacity as a function of truncated visit history
($T = 1, 2, 3, \ldots$).

\subsection{Latent Space Interpretability}
UMAP cluster separation and Integrated-Gradients (\texttt{captum}) region
contribution maps on the frozen GAAE/VGAE encoders.

\section{Discussion \& Limitations}
\begin{itemize}
  \item Inductive bias vs.\ model capacity in sample-constrained
    neuroimaging: a lightweight recurrent model over frozen graph embeddings
    matches or beats a much larger end-to-end transformer on this cohort
    size.
  \item Threats to validity: censoring/dropout patterns, small held-out test
    set, GPU nondeterminism as a general risk for ``stability'' claims about
    baselines (see Section~5) --- same-seed replicate runs, not just
    multi-seed spread, are needed to verify reproducibility claims.
  \item Clinical translation: current results are cross-sectional-adjacent
    (conversion by end of follow-up), not yet a calibrated time-to-event
    estimate.
\end{itemize}

\section{Conclusion \& Next Steps}

Three results are robust to the caveats above and are what the thesis will be built on:
(i) the published Brain-TokenGT baseline is not reproducible at a fixed seed, and its
multi-seed error bar misattributes run-to-run noise as seed sensitivity, while our pipeline
re-runs byte-identically; (ii) removing the graph encoder entirely does not degrade
performance, which extends a published cross-sectional, large-$N$ finding into the
longitudinal, small-$N$ regime; and (iii) $\Delta t$ conditioning cannot be evidenced on a
protocol-driven cohort whose inter-visit intervals are 90\% identical.

\paragraph{Next steps, in priority order.}
\begin{enumerate}
  \item \textbf{ADNI external validation of the \texttt{none} and
    \texttt{pretrained\_frozen} arms.} This is the highest-value open item by a wide
    margin: it converts a single-cohort negative result into a cross-cohort one, and it is
    the \emph{only} setting in which the $\Delta t$ mechanism can be tested, since ADNI
    encodes actual elapsed days rather than scheduled protocol months. The blocker is one
    wiring change --- the GE-LSTM dataset still imports the DELCODE-only month parser ---
    gated on a byte-equality check that DELCODE results are unchanged.
  \item \textbf{Rewrite the comparison notebook to glob all runs per experiment id} and
    drop both invalid paired tests, reporting distributions and an effect size instead.
  \item \textbf{Repeat sweep for the fine-tuned arm}, plus an \texttt{encoder\_lr\_scale}
    arm, before any claim about fine-tuning is made.
  \item Finalise the visit-horizon analysis in both cohort constructions, and the
    interpretability maps framed descriptively rather than as biomarker discovery.
\end{enumerate}

\paragraph{Explicitly not being done.} Chasing full CUDA determinism for Brain-TokenGT.
\texttt{torch.use\_deterministic\_algorithms(True)} would raise on operations that then
need replacing inside a third-party port, with no bounded time estimate. The 19-run
distribution is both cheaper and a stronger result than a repaired point estimate would be.


\begin{thebibliography}{9}
\bibitem{placeholder} Placeholder --- populate once source list is final.
\end{thebibliography}

\end{document}
